\begin{publicationbox}{\resizebox{0.5cm}{!}{\faFileText}}
    \textbf{F. Adesanya}, K. C. Silva, V. V. Graciano Neto, and I. David. ``Systems of Twinned Systems: A Systematic Review''. In: Software and Systems Modeling. 2026

    \tcblower
    
    The contents of this chapter have been accepted for publication in the International Journal on Software and Systems Modeling (SoSyM).
\end{publicationbox}

Systems of twinned systems (SoTS) arise at the intersection of systems of systems (SoS) and digital twins (DTs). Research on SoS focuses on coordinating autonomous, heterogeneous systems to achieve collective goals~\cite{maier1998architecting}, while research on DTs emphasizes continuously synchronized digital representations of physical assets that support monitoring, prediction, and control~\cite{kritzinger2018digital, tao2018digital}. As these research directions converge, an increasing body of work describes systems in which multiple digitally twinned systems interact, coordinate, and evolve under SoS principles.

\paragraph{Evidence of SoTS in existing research.}
This convergence is already evident across several cyber--physical domains. Smart-city and intelligent mobility systems provide prominent examples, where multiple digital twins are integrated to support coordination across large-scale, distributed systems. City-level digital twin architectures aggregate infrastructure-, vehicle-, and service-level twins to enable monitoring, analysis, and optimization across interconnected urban subsystems~\cite{tao2019digital, schroeder2021digital, gill2024toward}. In these environments, constituent systems typically retain operational autonomy and may dynamically join or leave coordinated activities.

Several studies further emphasize self-organization and self-optimization mechanisms, in which local decisions informed by digital twins contribute to emergent system-level behaviour~\cite{esterle2021digital, ghanbarifard2023digital}. Autonomous vehicle platooning in urban traffic illustrates this pattern, with coordination emerging through loosely coupled interactions among multiple digital twins rather than centralized control~\cite{borth2019digital, olsson2023systems}. Collectively, these studies reveal recurring structural and behavioural patterns consistent with SoTS.

\paragraph{Characterization SoTS.}
Taken together, this body of work indicates the emergence of a class of systems in which digitally twinned systems coordinate under SoS principles. While existing studies address specific applications, architectures, and coordination mechanisms, their broader implications for system structure and behaviour are not immediately evident when considered in isolation. Characterizing SoTS therefore requires a systematic synthesis of the literature to infer the common properties that define this class of systems and distinguish it from conventional SoS or standalone digital twin deployments.

\paragraph{Review objective and method.}
To characterize SoTS, this chapter presents a systematic literature review conducted according to the study design described in \secref{sec:literature-review-study-design}. The review applies explicit inclusion, exclusion, and data extraction criteria to analyze how existing research combines DTs with SoS principles. The chapter then introduces a classification framework for SoTS (\secref{sec:literature-review-classification}). This framework is used to identify and compare recurring architectural patterns and coordination strategies across the literature. The results of the review are presented in \secref{sec:literature-review-results} and summarized through a set of key takeaways that ground the problem addressed in subsequent chapters.